% ══════════════════════════════════════════════════════════════════════════
% main.tex — Reporte del Proyecto Parcial 1
% Unidad combinacional de 2 bits (MUX, decoder, display de 7 segmentos)
% ══════════════════════════════════════════════════════════════════════════
% Compile:  latexmk -pdf -auxdir=build main.tex
% ══════════════════════════════════════════════════════════════════════════

\documentclass[a4paper,12pt]{article}

% ══════════════════════════════════════════════
% config.tex — Configuración compartida
% ══════════════════════════════════════════════

\usepackage[utf8]{inputenc}
% ── Idioma ──
% Requiere texlive-langspanish (spanish.ldf + hyph-es.tex). Sin ese paquete
% \usepackage[spanish,es-nodecimaldot]{babel} aborta con "Unknown option 'spanish'".
% es-nodecimaldot: sin esta opcion babel-espanol cambia 0.7856 por 0,7856
\usepackage[spanish,es-nodecimaldot]{babel}
% babel-espanol mete un espacio fino antes de '%' incluso en modo
% math, lo que rompe \,74\,\% con "Incompatible glue units".
\spanishplainpercent
\usepackage{csquotes}
\usepackage[T1]{fontenc}
\usepackage[left=6cm,right=1.5cm,top=2.5cm,bottom=2.5cm]{geometry}
\usepackage[colorlinks=true, linkcolor=black, citecolor=black, urlcolor=blue]{hyperref}
\usepackage{amsmath}
\usepackage{amssymb}
\usepackage{wasysym}
\usepackage{booktabs}
\usepackage[version=4]{mhchem}
\usepackage{graphicx}
\usepackage{xcolor}
\usepackage{float}
\usepackage{enumitem}
\usepackage[export]{adjustbox}
\usepackage{epigraph}
\usepackage{tikz}
% Imprescindible con babel español: los caracteres < y > se vuelven activos y
% rompen la sintaxis de claves de TikZ (p. ej. ">=stealth"), provocando
% "Argument of \language@active@arg> has an extra }".
\usetikzlibrary{babel,arrows.meta}
\usepackage{helvet}
\usepackage{mathptmx}
\usepackage{array}
\usepackage[table]{xcolor}
\usepackage{caption}
\usepackage{mathtools}

% ── Encabezados y pies de página ──
\usepackage{fancyhdr}

% ── Marca de agua vertical en el margen ──
\usepackage{draftwatermark}
\SetWatermarkColor{blue}
\SetWatermarkAngle{90}
\SetWatermarkFontSize{9pt}
\SetWatermarkHorCenter{0.03\paperwidth}
\SetWatermarkText{Universidad Nacional Aut\'{o}noma de M\'{e}xico, Facultad de Ingenier\'{i}a}

% ── Bibliografía ──
\usepackage[
  backend=biber,
  style=authoryear,
  sorting=nyt,
  maxcitenames=2,
  uniquelist=false,
  hyperref=true
]{biblatex}
\addbibresource{config/references.bib}

% ── Unificar hipervínculos en citas (author+year en un solo link) ──
\DeclareCiteCommand{\parencite}[\mkbibparens]
  {\usebibmacro{prenote}}
  {\usebibmacro{citeindex}%
   \printtext[bibhyperref]{\usebibmacro{cite}}}
  {\multicitedelim}
  {\usebibmacro{postnote}}

\DeclareCiteCommand{\textcite}
  {\usebibmacro{prenote}}
  {\usebibmacro{citeindex}%
   \printtext[bibhyperref]{\usebibmacro{textcite}}}
  {\multicitedelim}
  {\usebibmacro{textcite:postnote}}

% ══════════════════════════════════════════════
%  NOTAS DE COMPILACIÓN (leer antes de tocar)
% ══════════════════════════════════════════════
%
% Dependencias del sistema (Arch):
%   sudo pacman -S texlive-langspanish biber
%
% 1) Backend: se usa biber, NO bibtex. Con backend=bibtex, biblatex 3.21 emite
%    "Using fall-back bibtex backend" en cada compilación y ese aviso no se
%    puede silenciar desde el preámbulo. Migrado a biber el 2026-10-05.
%
% 2) biber NO está en /usr/bin en Arch: el paquete lo instala como
%    /usr/bin/vendor_perl/biber. Si latexmk se queja de que no encuentra
%    biber, exporta el path:
%      export PATH="$PATH:/usr/bin/vendor_perl"
%    (ya está añadido a ~/.bashrc).
%
% 3) latexmk detecta biber solo, así que el comando normal es:
%      latexmk -pdf -auxdir=build main.tex
%    Con -bibtex también funciona, pero es innecesario y confuso: el flag
%    venía de cuando el backend era bibtex y latexmk -pdf a secas NO lo
%    ejecutaba, dejando las citas indefinidas sin avisar.
%
%    -auxdir=build mete TODOS los artefactos (.aux .bbl .bcf .blg .log .fls
%    .fdb_latexmk .out .run.xml) en build/ y deja main.pdf en el nivel
%    superior. Por eso el grep de errores va contra build/main.log, no
%    contra main.log:
%      grep -cE '^!' build/main.log
%    Ojo: -auxdir requiere latexmk >= 4.70 (aqui 4.87). Y si anades
%    % ══════════════════════════════════════════════
% config.tex — Configuración compartida
% ══════════════════════════════════════════════

\usepackage[utf8]{inputenc}
% ── Idioma ──
% Requiere texlive-langspanish (spanish.ldf + hyph-es.tex). Sin ese paquete
% \usepackage[spanish,es-nodecimaldot]{babel} aborta con "Unknown option 'spanish'".
% es-nodecimaldot: sin esta opcion babel-espanol cambia 0.7856 por 0,7856
\usepackage[spanish,es-nodecimaldot]{babel}
% babel-espanol mete un espacio fino antes de '%' incluso en modo
% math, lo que rompe \,74\,\% con "Incompatible glue units".
\spanishplainpercent
\usepackage{csquotes}
\usepackage[T1]{fontenc}
\usepackage[left=6cm,right=1.5cm,top=2.5cm,bottom=2.5cm]{geometry}
\usepackage[colorlinks=true, linkcolor=black, citecolor=black, urlcolor=blue]{hyperref}
\usepackage{amsmath}
\usepackage{amssymb}
\usepackage{wasysym}
\usepackage{booktabs}
\usepackage[version=4]{mhchem}
\usepackage{graphicx}
\usepackage{xcolor}
\usepackage{float}
\usepackage{enumitem}
\usepackage[export]{adjustbox}
\usepackage{epigraph}
\usepackage{tikz}
% Imprescindible con babel español: los caracteres < y > se vuelven activos y
% rompen la sintaxis de claves de TikZ (p. ej. ">=stealth"), provocando
% "Argument of \language@active@arg> has an extra }".
\usetikzlibrary{babel,arrows.meta}
\usepackage{helvet}
\usepackage{mathptmx}
\usepackage{array}
\usepackage[table]{xcolor}
\usepackage{caption}
\usepackage{mathtools}

% ── Encabezados y pies de página ──
\usepackage{fancyhdr}

% ── Marca de agua vertical en el margen ──
\usepackage{draftwatermark}
\SetWatermarkColor{blue}
\SetWatermarkAngle{90}
\SetWatermarkFontSize{9pt}
\SetWatermarkHorCenter{0.03\paperwidth}
\SetWatermarkText{Universidad Nacional Aut\'{o}noma de M\'{e}xico, Facultad de Ingenier\'{i}a}

% ── Bibliografía ──
\usepackage[
  backend=biber,
  style=authoryear,
  sorting=nyt,
  maxcitenames=2,
  uniquelist=false,
  hyperref=true
]{biblatex}
\addbibresource{config/references.bib}

% ── Unificar hipervínculos en citas (author+year en un solo link) ──
\DeclareCiteCommand{\parencite}[\mkbibparens]
  {\usebibmacro{prenote}}
  {\usebibmacro{citeindex}%
   \printtext[bibhyperref]{\usebibmacro{cite}}}
  {\multicitedelim}
  {\usebibmacro{postnote}}

\DeclareCiteCommand{\textcite}
  {\usebibmacro{prenote}}
  {\usebibmacro{citeindex}%
   \printtext[bibhyperref]{\usebibmacro{textcite}}}
  {\multicitedelim}
  {\usebibmacro{textcite:postnote}}

% ══════════════════════════════════════════════
%  NOTAS DE COMPILACIÓN (leer antes de tocar)
% ══════════════════════════════════════════════
%
% Dependencias del sistema (Arch):
%   sudo pacman -S texlive-langspanish biber
%
% 1) Backend: se usa biber, NO bibtex. Con backend=bibtex, biblatex 3.21 emite
%    "Using fall-back bibtex backend" en cada compilación y ese aviso no se
%    puede silenciar desde el preámbulo. Migrado a biber el 2026-10-05.
%
% 2) biber NO está en /usr/bin en Arch: el paquete lo instala como
%    /usr/bin/vendor_perl/biber. Si latexmk se queja de que no encuentra
%    biber, exporta el path:
%      export PATH="$PATH:/usr/bin/vendor_perl"
%    (ya está añadido a ~/.bashrc).
%
% 3) latexmk detecta biber solo, así que el comando normal es:
%      latexmk -pdf -auxdir=build main.tex
%    Con -bibtex también funciona, pero es innecesario y confuso: el flag
%    venía de cuando el backend era bibtex y latexmk -pdf a secas NO lo
%    ejecutaba, dejando las citas indefinidas sin avisar.
%
%    -auxdir=build mete TODOS los artefactos (.aux .bbl .bcf .blg .log .fls
%    .fdb_latexmk .out .run.xml) en build/ y deja main.pdf en el nivel
%    superior. Por eso el grep de errores va contra build/main.log, no
%    contra main.log:
%      grep -cE '^!' build/main.log
%    Ojo: -auxdir requiere latexmk >= 4.70 (aqui 4.87). Y si anades
%    % ══════════════════════════════════════════════
% config.tex — Configuración compartida
% ══════════════════════════════════════════════

\usepackage[utf8]{inputenc}
% ── Idioma ──
% Requiere texlive-langspanish (spanish.ldf + hyph-es.tex). Sin ese paquete
% \usepackage[spanish,es-nodecimaldot]{babel} aborta con "Unknown option 'spanish'".
% es-nodecimaldot: sin esta opcion babel-espanol cambia 0.7856 por 0,7856
\usepackage[spanish,es-nodecimaldot]{babel}
% babel-espanol mete un espacio fino antes de '%' incluso en modo
% math, lo que rompe \,74\,\% con "Incompatible glue units".
\spanishplainpercent
\usepackage{csquotes}
\usepackage[T1]{fontenc}
\usepackage[left=6cm,right=1.5cm,top=2.5cm,bottom=2.5cm]{geometry}
\usepackage[colorlinks=true, linkcolor=black, citecolor=black, urlcolor=blue]{hyperref}
\usepackage{amsmath}
\usepackage{amssymb}
\usepackage{wasysym}
\usepackage{booktabs}
\usepackage[version=4]{mhchem}
\usepackage{graphicx}
\usepackage{xcolor}
\usepackage{float}
\usepackage{enumitem}
\usepackage[export]{adjustbox}
\usepackage{epigraph}
\usepackage{tikz}
% Imprescindible con babel español: los caracteres < y > se vuelven activos y
% rompen la sintaxis de claves de TikZ (p. ej. ">=stealth"), provocando
% "Argument of \language@active@arg> has an extra }".
\usetikzlibrary{babel,arrows.meta}
\usepackage{helvet}
\usepackage{mathptmx}
\usepackage{array}
\usepackage[table]{xcolor}
\usepackage{caption}
\usepackage{mathtools}

% ── Encabezados y pies de página ──
\usepackage{fancyhdr}

% ── Marca de agua vertical en el margen ──
\usepackage{draftwatermark}
\SetWatermarkColor{blue}
\SetWatermarkAngle{90}
\SetWatermarkFontSize{9pt}
\SetWatermarkHorCenter{0.03\paperwidth}
\SetWatermarkText{Universidad Nacional Aut\'{o}noma de M\'{e}xico, Facultad de Ingenier\'{i}a}

% ── Bibliografía ──
\usepackage[
  backend=biber,
  style=authoryear,
  sorting=nyt,
  maxcitenames=2,
  uniquelist=false,
  hyperref=true
]{biblatex}
\addbibresource{config/references.bib}

% ── Unificar hipervínculos en citas (author+year en un solo link) ──
\DeclareCiteCommand{\parencite}[\mkbibparens]
  {\usebibmacro{prenote}}
  {\usebibmacro{citeindex}%
   \printtext[bibhyperref]{\usebibmacro{cite}}}
  {\multicitedelim}
  {\usebibmacro{postnote}}

\DeclareCiteCommand{\textcite}
  {\usebibmacro{prenote}}
  {\usebibmacro{citeindex}%
   \printtext[bibhyperref]{\usebibmacro{textcite}}}
  {\multicitedelim}
  {\usebibmacro{textcite:postnote}}

% ══════════════════════════════════════════════
%  NOTAS DE COMPILACIÓN (leer antes de tocar)
% ══════════════════════════════════════════════
%
% Dependencias del sistema (Arch):
%   sudo pacman -S texlive-langspanish biber
%
% 1) Backend: se usa biber, NO bibtex. Con backend=bibtex, biblatex 3.21 emite
%    "Using fall-back bibtex backend" en cada compilación y ese aviso no se
%    puede silenciar desde el preámbulo. Migrado a biber el 2026-10-05.
%
% 2) biber NO está en /usr/bin en Arch: el paquete lo instala como
%    /usr/bin/vendor_perl/biber. Si latexmk se queja de que no encuentra
%    biber, exporta el path:
%      export PATH="$PATH:/usr/bin/vendor_perl"
%    (ya está añadido a ~/.bashrc).
%
% 3) latexmk detecta biber solo, así que el comando normal es:
%      latexmk -pdf -auxdir=build main.tex
%    Con -bibtex también funciona, pero es innecesario y confuso: el flag
%    venía de cuando el backend era bibtex y latexmk -pdf a secas NO lo
%    ejecutaba, dejando las citas indefinidas sin avisar.
%
%    -auxdir=build mete TODOS los artefactos (.aux .bbl .bcf .blg .log .fls
%    .fdb_latexmk .out .run.xml) en build/ y deja main.pdf en el nivel
%    superior. Por eso el grep de errores va contra build/main.log, no
%    contra main.log:
%      grep -cE '^!' build/main.log
%    Ojo: -auxdir requiere latexmk >= 4.70 (aqui 4.87). Y si anades
%    \input{config/config} en vez de \input{config}, y \addbibresource
%    con ruta completa desde la raiz del proyecto, porque al estar el
%    .bcf en build/ las rutas relativas se resuelven distinto.
%
% 4) Tipos de entrada: usa @online para páginas web. @webpage no existe en
%    biblatex y produce "No driver for 'webpage'" en cada compilación.
%    Los .bib de la vault se migraron de @webpage a @online.
 en vez de \input{config}, y \addbibresource
%    con ruta completa desde la raiz del proyecto, porque al estar el
%    .bcf en build/ las rutas relativas se resuelven distinto.
%
% 4) Tipos de entrada: usa @online para páginas web. @webpage no existe en
%    biblatex y produce "No driver for 'webpage'" en cada compilación.
%    Los .bib de la vault se migraron de @webpage a @online.
 en vez de \input{config}, y \addbibresource
%    con ruta completa desde la raiz del proyecto, porque al estar el
%    .bcf en build/ las rutas relativas se resuelven distinto.
%
% 4) Tipos de entrada: usa @online para páginas web. @webpage no existe en
%    biblatex y produce "No driver for 'webpage'" en cada compilación.
%    Los .bib de la vault se migraron de @webpage a @online.


% ──────────────────────────────────────────────────────────────────────
%  METADATA
% ──────────────────────────────────────────────────────────────────────
\newcommand{\materia}{Dise\~{n}o Digital Moderno}
\newcommand{\tarea}{Proyecto Parcial 1: Unidad combinacional de 2 bits con MUX, decoder y display de 7 segmentos}
\newcommand{\grupo}{5}
\newcommand{\semestre}{2027-1}
\newcommand{\maestro}{Dr. Oscar Francisco Fuentes Casarrubias}
\newcommand{\fechaentrega}{29 de octubre de 2026}

\DeclareUnicodeCharacter{2212}{-}

\hypersetup{
  pdfauthor={Hernandez Rosas Jafet Daniel},
  pdftitle={\materia - \tarea},
  pdflang={es}
}

\emergencystretch=2em

% ──────────────────────────────────────────────────────────────────────
%  DOCUMENT
% ──────────────────────────────────────────────────────────────────────
\begin{document}
\thispagestyle{empty}
\SetWatermarkText{}

% ══════════════════════════════════════════════════════════════════════
%  PORTADA
% ══════════════════════════════════════════════════════════════════════

\begin{tikzpicture}[remember picture, overlay]
    \node[anchor=north] (unam) at ([xshift=2.8cm, yshift=-2cm]current page.north west) {
        \includegraphics[width=3.5cm]{images/UNAM.png}
    };
    \node[anchor=south] (fi) at ([xshift=2.8cm, yshift=2cm]current page.south west) {
        \includegraphics[width=3.5cm]{images/FI.png}
    };
    \draw[double, double distance=12pt, line width=2.5pt]
        ([yshift=-0.5cm]unam.south) -- ([yshift=0.5cm]fi.north);
\end{tikzpicture}

\begin{center}
    \vspace{-40px}
    \hspace*{0.5cm}
    \begin{minipage}{\dimexpr\linewidth-0.5cm-3pt\relax}
        \centering
        \textbf{\large UNIVERSIDAD NACIONAL AUT\'{O}NOMA DE \\ M\'{E}XICO}\\[0.2cm]
        \rule{\linewidth}{2pt} \\[0.5cm]
        \vspace{-12px}
        \textbf{\large FACULTAD DE INGENIER\'{I}A}\\[1.4cm]
        \textbf{\large{\materia}}\\[1.5cm]
    \end{minipage}
\end{center}
\vspace{10px}

\noindent
{\large\raggedright
\textbf{Tarea:} \tarea\\[0.4cm]
\textbf{Profesor:} \maestro\\[0.2cm]
\textbf{Semestre:} \semestre\\[0.6cm]
\textbf{Integrantes:} Duarte Puertas Manuel \\
Hernandez Rosas Jafet Daniel \\
Reyes Garc\'{i}a Miguel Angel \\
Torres Rodriguez Lizeth Danae\\[0.2cm]
\textbf{Grupo:} \grupo\\[0.2cm]
\textbf{Fecha de entrega:} \fechaentrega\par
}

% ══════════════════════════════════════════════════════════════════════
%  CONTENIDO
% ══════════════════════════════════════════════════════════════════════
\newpage
\setcounter{page}{1}
\newgeometry{left=2.5cm, right=2.5cm, top=2.5cm, bottom=2.5cm}

\pagestyle{fancy}
\fancyhf{}
\setlength{\headheight}{15pt}
\fancyhead[L]{\materia}
\fancyhead[R]{Proyecto Parcial 1}
\fancyfoot[C]{\thepage}
\renewcommand{\headrulewidth}{0.4pt}

\SetWatermarkText{Universidad Nacional Aut\'{o}noma de M\'{e}xico, Facultad de Ingenier\'{i}a}

% ─────────────────────────────────────────────────────────────────────
\section{Introducci\'{o}n}
En este reporte se documenta el dise\~{n}o, montaje y verificaci\'{o}n
de una unidad aritm\'{e}tico-l\'{o}gica m\'{i}nima de dos bits, construida
\'{u}nicamente con circuitos integrados discretos de la familia TTL LS.
El sistema recibe dos operandos \texttt{A = A1A0} y \texttt{B = B1B0}
de dos bits, una se\~{n}al de selecci\'{o}n \texttt{S = S1S0} y una
se\~{n}al de habilitaci\'{o}n \texttt{ENABLE}, y produce un resultado
en el rango $\{0,1,\ldots,6\}$ que se visualiza en un display de
\'{a}nodo com\'{u}n de siete segmentos.

Las cuatro operaciones implementadas son \texttt{SUM} (suma binaria sin
signo), \texttt{XOR} y \texttt{AND} bit a bit, y \texttt{MAX} (mayor
entre \texttt{A} y \texttt{B}). La selecci\'{o}n se hace con un
multiplexor 74LS153, mientras que un decoder 74LS139 activa un LED
indicador de la operaci\'{o}n seleccionada. El resultado final pasa
por un driver BCD a 7 segmentos 74LS47.

% ─────────────────────────────────────────────────────────────────────
\section{Objetivo}
Dise\~{n}ar, cablear y verificar una unidad l\'{o}gica combinacional de
dos bits que, a partir de dos operandos \texttt{A} y \texttt{B}, una
selecci\'{o}n \texttt{S} y una habilitaci\'{o}n \texttt{ENABLE},
entregue el resultado de la operaci\'{o}n elegida en un display de
siete segmentos, usando exclusivamente compuertas, multiplexores,
decoders y drivers TTL.

% ─────────────────────────────────────────────────────────────────────
\section{Arquitectura general}
El sistema se parti\'{o} en dos protoboards de 830 puntos seg\'{u}n la
distribuci\'{o}n f\'{i}sica de la Figura~\ref{fig:topologia}. La
\textit{Placa 1} contiene las entradas y los bloques que calculan las
cuatro operaciones. La \textit{Placa 2} contiene la etapa de
selecci\'{o}n (MUX), el decoder, los LEDs indicadores, el driver BCD y
el display.

\begin{figure}[H]
\centering
\begin{tikzpicture}[
    every node/.style={font=\small},
    blk/.style  ={draw, rounded corners, align=center, fill=blue!5,
                  minimum width=2.6cm, minimum height=1.0cm},
    sw/.style   ={draw, fill=yellow!15, align=center, minimum width=1.6cm,
                  minimum height=0.5cm, font=\scriptsize},
    led/.style  ={circle, draw, minimum size=0.4cm, fill=red!10, inner sep=0pt},
    arr/.style  ={->, thick},
    lbl/.style  ={font=\scriptsize, fill=white, inner sep=1pt,
                  fill opacity=0.85, text opacity=1}
]

%==========================================================================
% Bloques. Coordenadas explcitas; cada bloque usa minimum width = X cm.
% X de borde izquierdo = center.x - X/2, X de borde derecho = center.x + X/2.
%
% sw1, sw2:      center (0.8, 4.5) y (0.8, 2.5).  min width 1.6 -> bordes en x=0 y 1.6.
% ops:           center (4.5, 3.5).                min width 2.6 -> bordes 3.2 y 5.8.
% mux, dec, drv: center x = 8.5.                   min width 2.6 -> bordes 7.2 y 9.8.
%                 y = 4.5, 2.5, 0.5.
% dsp:           center (12.0, 0.5).              min width 2.6 -> bordes 10.7 y 13.3.
% LEDs:          center (12.0, 3.5-i*0.5).
%==========================================================================

\node[sw]  (sw1) at (0.8, 4.5) {SW1};
\node[sw]  (sw2) at (0.8, 2.5) {SW2};
\node[blk] (ops) at (4.5, 3.5)
    {\textbf{Operaciones}\\\scriptsize SUM, XOR, AND, MAX};
\node[blk] (mux) at (8.5, 4.5) {MUX\\\scriptsize 74LS153};
\node[blk] (dec) at (8.5, 2.5) {Decoder\\\scriptsize 74LS139};
\node[blk] (drv) at (8.5, 0.5) {Driver\\\scriptsize 74LS47};
\node[blk, minimum width=2.6cm] (dsp) at (12.0, 0.5)
    {Display 7-seg\\\scriptsize \'{a}nodo};

\foreach \i/\lab in {0/MAX, 1/AND, 2/XOR, 3/SUM}{
    \node[led] (L\i) at (12.0, 3.5-\i*0.5) {};
    \node[anchor=west, font=\scriptsize] at (12.3, 3.5-\i*0.5)
        {LED\i: \lab};
}

% Flechas con coordenadas LITERALES, no `|-`.
% sw1 (east edge x=1.6, y=4.5) -> ops (west edge x=3.2, y=3.5).
%   Cable: sal horizontal desde sw1, luego baja vertical al oeste de ops.
\draw[thick] (1.6, 4.5) -- (3.2, 4.5);
\draw[arr]  (3.2, 4.5) -- (3.2, 3.5);  % punta de flecha al lado oeste de ops
\node[lbl, anchor=south, font=\scriptsize] at (2.4, 4.6) {$A_1,A_0,B_1,B_0$};

% ops (east edge x=5.8, y=3.5) -> mux (west edge x=7.2, y=4.5).
%   Subir desde y=3.5 hasta y=4.5, luego horizontal, luego entrar a mux.
\draw[thick] (5.8, 3.5) -- (5.8, 4.5);   % sube vertical
\draw[arr]  (5.8, 4.5) -- (7.2, 4.5);
\node[lbl, anchor=south, font=\scriptsize] at (6.5, 4.6)
    {$\Sigma_i, X_i, N_i, M_i$};

% sw2 (east edge x=1.6, y=2.5) -> dec (west edge x=7.2, y=2.5).
\draw[arr] (1.6, 2.5) -- (7.2, 2.5);
\node[lbl, anchor=south, font=\scriptsize] at (4.4, 2.6)
    {$S_1, S_0$};

% mux -> drv (flecha vertical via dec)
\draw[arr] (mux.south) -- (dec.north)
    node[lbl, midway, right]{$R_i$};
\draw[arr] (dec.south) -- (drv.north);

% /EN: indicador pequeno entre mux y dec (no es cable real, solo anotacion)
\node[lbl, anchor=west, font=\scriptsize] at (8.6, 3.0) {$\overline{\mathrm{EN}}$};
\draw[thick, dashed] (7.2, 3.5) -- (7.2, 3.0) -- (8.4, 3.0) -- (8.4, 3.5);

% driver -> display
\draw[arr] (9.8, 0.5) -- (10.7, 0.5);

% decoder -> 4 LEDs (lineas diagonales cortas)
\draw[arr] (9.8, 2.5) -- (12.0, 3.5);
\draw[arr] (9.8, 2.5) -- (12.0, 3.0);
\draw[arr] (9.8, 2.5) -- (12.0, 2.5);
\draw[arr] (9.8, 2.5) -- (12.0, 2.0);

\end{tikzpicture}
\caption{Diagrama de bloques funcional del sistema. Las flechas
horizontales con etiquetas representan los buses principales
(\texttt{A,B}, \texttt{$\Sigma_i,X_i,N_i,M_i$}, \texttt{$S_1,S_0$}). La
flecha etiquetada \texttt{$R_i$} indica el flujo del resultado
seleccionado por el MUX hacia el driver; las cuatro l\'{i}neas que
salen del decoder energizan los LEDs indicadores.}
\label{fig:topologia}
\end{figure}

Cada uno de los bloques \texttt{SUM}, \texttt{XOR}, \texttt{AND} y
\texttt{MAX} entrega sus dos bits de salida al MUX de la Placa 2. All\'{i}
el 74LS153, gobernado por \texttt{S1\,S0} y \texttt{EN}, elige una de las
cuatro operaciones y la env\'{i}a al driver BCD 74LS47, que a su vez
maneja el display. En paralelo, el 74LS139 decodifica \texttt{S1\,S0} en
cuatro salidas activas en bajo que encienden un LED por operaci\'{o}n
cuando \texttt{EN = 1}.

% ─────────────────────────────────────────────────────────────────────
\section{Tabla de especificaciones}

\begin{table}[H]
\centering
\small
\setlength{\tabcolsep}{4pt}
\begin{tabular}{p{3.2cm}p{12.3cm}}
\toprule
\textbf{Nombre del sistema}     & Unidad combinacional de 2 bits (ALU\_2b v1).\\
\textbf{Entradas}               & \texttt{A1, A0, B1, B0} (operandos, dos DIP switches);
                                   \texttt{S1, S0} (selecci\'{o}n de operaci\'{o}n);
                                   \texttt{ENABLE} (habilitaci\'{o}n).\\
\textbf{Salidas}                & \texttt{R2, R1, R0} en BCD al driver 74LS47;
                                   cuatro LEDs indicadores de la operaci\'{o}n activa;
                                   display de 7 segmentos de \'{a}nodo com\'{u}n.\\
\textbf{Rango de operaci\'{o}n}& \texttt{A}, \texttt{B} $\in \{0,1,2,3\}$ sin signo.
                                   Resultado $\in \{0,1,\ldots,6\}$.\\
\textbf{Alimentaci\'{o}n}        & $+5\,\mathrm{V}$ regulados; tierra com\'{u}n;
                                   capacitores de desacoplo de $100\,\mathrm{nF}$
                                   junto a cada CI.\\
\textbf{CIs utilizados}         & 74LS283, 74LS86, 74LS08, 74LS85, 74LS157,
                                   74LS04, 74LS139, 74LS153, 74LS47.\\
\textbf{Pasivos}                & 4 DIP switches (8 + 4 posiciones),
                                   4 LEDs con resistencias de $470\,\Omega$,
                                   display de \'{a}nodo com\'{u}n con resistencias
                                   de $330\,\Omega$ por segmento, y 12
                                   resistencias de pull-down de $10\,\mathrm{k}\Omega$.\\
\bottomrule
\end{tabular}
\caption{Hoja de especificaciones (versi\'{o}n resumida; la versi\'{o}n
a una p\'{a}gina se entrega por separado en
\texttt{Especificaciones/}).}
\label{tab:especificaciones}
\end{table}

\subsection{Codificaci\'{o}n de operaciones \texttt{S1S0}}
El circuito cableado asign\'{o} las operaciones a las entradas del MUX
74LS153 en el siguiente orden (importante: \textbf{el orden difiere del
enunciado}, ver Secci\'{o}n~\ref{sec:observaciones}):

\begin{table}[H]
\centering
\begin{tabular}{cc|l}
\toprule
\texttt{S1} & \texttt{S0} & Operaci\'{o}n seleccionada \\
\midrule
0 & 0 & \textbf{MAX} (mayor entre $A$ y $B$) \\
0 & 1 & \textbf{AND} bit a bit \\
1 & 0 & \textbf{XOR} bit a bit \\
1 & 1 & \textbf{SUM} ($A + B$) \\
\bottomrule
\end{tabular}
\caption{Mapeo cableado entre \texttt{S1S0} y la operaci\'{o}n que llega
al display. Esta tabla refleja lo armado en protoboard; el enunciado
original propon\'{i}a el orden inverso.}
\label{tab:smux}
\end{table}

\subsection{Comportamiento de \texttt{ENABLE}}
La se\~{n}al \texttt{ENABLE} se invierte con un 74LS04 y la se\~{n}al
negada \texttt{$\overline{\text{EN}}$} se conecta a las entradas de
habilitaci\'{o}n activas en bajo del MUX (74LS153, pin 1) y del decoder
(74LS139, pin 1).

\begin{itemize}[leftmargin=*, itemsep=2pt]
  \item \texttt{EN = 1}: \texttt{$\overline{\text{EN}} = 0$}; MUX y
        decoder habilitados. El display muestra la operaci\'{o}n
        seleccionada y \'{u}nicamente un LED indicador queda encendido.
  \item \texttt{EN = 0}: \texttt{$\overline{\text{EN}} = 1$}; MUX y
        decoder forzados a su estado inactivo. El 74LS153 pone sus
        salidas en 0 y el 74LS139 pone las cuatro salidas en alto (l\'{o}gica
        inversa), por lo que \textbf{todos los LEDs se apagan y el display
        muestra 0}. Este comportamiento es sim\'{e}trico: el sistema est\'{a}
        definido en su estado deshabilitado y no quedan se\~{n}ales
        flotantes.
\end{itemize}

% ─────────────────────────────────────────────────────────────────────
\section{Desarrollo de las operaciones}
Cada operaci\'{o}n se realiza en la Placa 1 y entrega dos bits al MUX
de la Placa 2. En el caso de la suma se obtiene adem\'{a}s un bit
\textit{extra} \texttt{SUM2} que el MUX ignora; como $A, B \le 3$,
$A+B \le 6 < 8$, as\'{i} que \texttt{SUM2} nunca toma valor $1$ para
los datos de prueba del enunciado.

\subsection{Suma (\texttt{SUM}): 74LS283}
El 74LS283 es un sumador binario completo de cuatro bits con carry de
entrada y de salida. Lo usamos con las entradas $A_1A_0$ y $B_1B_0$ en
los pines $A_1, A_0, B_1, B_0$ (3, 5, 2, 6) y el acarreo de entrada
($C_\text{in}$, pin 8 en LS283) conectado a \texttt{GND}. Las tres
salidas bajas \texttt{$\Sigma_2\,\Sigma_1\,\Sigma_0$} se conectan al
MUX. La salida $\Sigma_3$ (pin 13) queda al aire en este dise\~{n}o.

\paragraph{Tabla de verdad (suma de dos n\'{u}meros de 2 bits).}
\begin{table}[H]
\centering
\small
\begin{tabular}{ccc|cc}
\toprule
$A_1$ & $A_0$ & $B_1\,B_0$ & $\Sigma_1\Sigma_0$ & $\Sigma_2$ \\
\midrule
0 & 0 & \textit{variable} & $B_1B_0$ & 0 \\
0 & 1 & \textit{variable} & $B_1B_0 \oplus 01$ & 0 \\
1 & 0 & \textit{variable} & $B_1B_0 \oplus 10$ & $B_1$ \\
1 & 1 & \textit{variable} & $\overline{B_1B_0}$ & $B_1 \lor B_0$ \\
\bottomrule
\end{tabular}
\caption{Salidas del 74LS283 vistas como funci\'{o}n de $A$ y $B$.}
\end{table}

\paragraph{Ecuaciones.}
\begin{align}
\Sigma_0 &= A_0 \oplus B_0 \oplus 0 = A_0 \oplus B_0, \label{eq:s0}\\
\Sigma_1 &= A_1 \oplus B_1 \oplus (A_0 \cdot B_0),              \label{eq:s1}\\
\Sigma_2 &= (A_1 \cdot B_1) \lor \bigl((A_1 \oplus B_1) \cdot (A_0 \cdot B_0)\bigr).
           \label{eq:s2}
\end{align}
La simplificaci\'{o}n es la habitual del sumador completo; en este
caso el acarreo de entrada es $0$, lo que reduce $\Sigma_0$ a un
simple XOR.

\paragraph{Implementaci\'{o}n.} Sin circuito adicional: el 74LS283
realiza las ecuaciones~\ref{eq:s0}--\ref{eq:s2} por su propia cuenta.

\subsection{XOR bit a bit: 74LS86}
\paragraph{Tabla de verdad.}
\begin{table}[H]
\centering
\begin{tabular}{cc|cc}
\toprule
$A_i$ & $B_i$ & $X_i = A_i \oplus B_i$ \\
\midrule
0 & 0 & 0 \\
0 & 1 & 1 \\
1 & 0 & 1 \\
1 & 1 & 0 \\
\bottomrule
\end{tabular}
\end{table}

\paragraph{Ecuaciones y simplificaci\'{o}n.}
$$X_0 = A_0\overline{B_0} \lor \overline{A_0}B_0, \qquad
  X_1 = A_1\overline{B_1} \lor \overline{A_1}B_1.$$
Ya en su forma m\'{i}nima (un solo XOR por bit); no requiere m\'{a}s
simplificaci\'{o}n.

\paragraph{Implementaci\'{o}n.} Dos de las cuatro compuertas XOR del
74LS86, una por bit.

\subsection{AND bit a bit: 74LS08}
\paragraph{Tabla y ecuaciones.}
$$N_0 = A_0 \cdot B_0, \qquad N_1 = A_1 \cdot B_1.$$
La funci\'{o}n AND ya es m\'{i}nima.

\paragraph{Implementaci\'{o}n.} Dos de las cuatro compuertas AND del
74LS08.

\subsection{MAX: 74LS85 + 74LS157}
\paragraph{Definici\'{o}n.} $\mathrm{MAX}(A,B) = A$ si $A \ge B$ y $B$
si $A < B$.

\paragraph{Tabla de verdad (A y B sin signo).}
\begin{table}[H]
\centering
\small
\begin{tabular}{cc|cccc}
\toprule
$A_1A_0$ & $B_1B_0$ & $M_1M_0$ & $A>B$ & $A=B$ & $A<B$ \\
\midrule
00 & 00 & 00 & 0 & 1 & 0 \\
00 & 01 & 01 & 0 & 0 & 1 \\
00 & 10 & 10 & 0 & 0 & 1 \\
00 & 11 & 11 & 0 & 0 & 1 \\
01 & 00 & 01 & 1 & 0 & 0 \\
01 & 01 & 01 & 0 & 1 & 0 \\
01 & 10 & 10 & 0 & 0 & 1 \\
01 & 11 & 11 & 0 & 0 & 1 \\
10 & 00 & 10 & 1 & 0 & 0 \\
10 & 01 & 10 & 1 & 0 & 0 \\
10 & 10 & 10 & 0 & 1 & 0 \\
10 & 11 & 11 & 0 & 0 & 1 \\
11 & 00 & 11 & 1 & 0 & 0 \\
11 & 01 & 11 & 1 & 0 & 0 \\
11 & 10 & 11 & 1 & 0 & 0 \\
11 & 11 & 11 & 0 & 1 & 0 \\
\bottomrule
\end{tabular}
\caption{Mapeo entrada $\rightarrow$ salida del bloque MAX. La columna
``>$/$=/$<$'' indica cu\'{a}l de las tres salidas activas en alto del
74LS85 se activa para cada par.}
\label{tab:max}
\end{table}

\paragraph{Ecuaciones.}
\begin{align}
M_0 &= \overline{GT}\cdot A_0 \lor GT\cdot B_0,\\
M_1 &= \overline{GT}\cdot A_1 \lor GT\cdot B_1,
\end{align}
donde $GT$ es la salida \texttt{PQA $>$ B} del 74LS85. En realidad la
salida es la suma l\'{o}gica (cableada) $GT \lor EQ$; como el
74LS157 solo tiene un pin de selecci\'{o}n por secci\'{o}n, conectamos
esa OR externa al select com\'{u}n del MUX.

\paragraph{Implementaci\'{o}n.} El 74LS85 produce las tres se\~{n}ales
\texttt{A$>$B}, \texttt{A$=$B} y \texttt{A$<$B}; con una OR externa de
dos entradas, \texttt{A$>$B $\lor$ A$=$B}, conectamos al pin de
selecci\'{o}n com\'{u}n del 74LS157. Las dos secciones del 157
seleccionan entre $A$ (cuando $\text{select}=0$) y $B$ (cuando
$\text{select}=1$). Para $A \le B$ el bloque entrega $A$; para $A > B$
entrega $B$, que es el mayor. Esto cubre los 16 casos de la
Tabla~\ref{tab:max}.

% ─────────────────────────────────────────────────────────────────────
\section{MUX y decoder}

\subsection{MUX 74LS153}
El 74LS153 contiene dos multiplexores 4-a-1 independientes que comparten
las se\~{n}ales \texttt{S1} y \texttt{S0} y tienen habilitaci\'{o}n
activa en bajo por separado. En este dise\~{n}o unimos ambas
habilitaciones para tratarlas como una sola y direccionamos una salida
por cada bit:

\begin{table}[H]
\centering
\small
\begin{tabular}{cc|cccc|cc}
\toprule
\texttt{S1} & \texttt{S0} &
\multicolumn{4}{c|}{Entradas (1C0, 1C1, 1C2, 1C3)} &
\multicolumn{2}{c}{Salida} \\
& & \footnotesize M(MAX) & \footnotesize N(AND) & \footnotesize X(XOR) & \footnotesize $\Sigma$(SUM) & \footnotesize R0 & \footnotesize R1 \\
\midrule
0 & 0 & 1C0 & --  & --   & --   & $M_0$ & $M_1$ \\
0 & 1 & --  & 1C1 & --   & --   & $N_0$ & $N_1$ \\
1 & 0 & --  & --  & 1C2  & --   & $X_0$ & $X_1$ \\
1 & 1 & --  & --  & --   & 1C3  & $\Sigma_0$ & $\Sigma_1$ \\
\bottomrule
\end{tabular}
\caption{Mapeo interno del 74LS153. La columna ``Entradas'' muestra, para
cada combinaci\'{o}n de \texttt{S1S0}, cu\'{a}l de los cuatro bloques de
operaci\'{o}n (filas superiores) queda conectada a la salida.}
\end{table}

% ─────────────────────────────────────────────────────────────────────

\subsection{Decoder 74LS139}
El 74LS139 es un decoder dual 2-a-4 con habilitaci\'{o}n activa en
bajo. Usamos una de sus mitades, con \texttt{$\overline{\text{EN}}$}
como habilitaci\'{o}n y \texttt{S1, S0} como entradas. Las cuatro
salidas activas en bajo atacan el \'{a}nodo de los LEDs indicadores (a
trav\'{e}s de las resistencias limitadoras de $470\,\Omega$). Con
\texttt{EN = 1} s\'{o}lo una salida est\'{a} en bajo, y por tanto s\'{o}lo
un LED enciende. Con \texttt{EN = 0} las cuatro salidas est\'{a}n en
alto y \textbf{ning\'{u}n} LED enciende.

\begin{table}[H]
\centering
\small
\begin{tabular}{cc|c|cccc}
\toprule
\texttt{EN} & \texttt{S1S0} & Habilitado & $\overline{Y_0}$ & $\overline{Y_1}$ & $\overline{Y_2}$ & $\overline{Y_3}$ \\
\midrule
0 & \textit{xx} & No & 1 & 1 & 1 & 1 \\
1 & 00 & S\'{i} & \textbf{0} & 1 & 1 & 1 \\
1 & 01 & S\'{i} & 1 & \textbf{0} & 1 & 1 \\
1 & 10 & S\'{i} & 1 & 1 & \textbf{0} & 1 \\
1 & 11 & S\'{i} & 1 & 1 & 1 & \textbf{0} \\
\bottomrule
\end{tabular}
\caption{Tabla funcional del 74LS139 con las cuatro se\~{n}ales
$\overline{Y_i}$ activas en bajo. Cada $\overline{Y_i}$ enciende el LED
correspondiente (MAX, AND, XOR, SUM en ese orden).}
\end{table}

% ─────────────────────────────────────────────────────────────────────
\section{Etapa de salida: driver BCD y display}

\paragraph{74LS47 y display de \'{a}nodo com\'{u}n.} El 74LS47 traduce
una entrada BCD ($\overline{R_3}\,R_2\,R_1\,R_0$ si contamos con la
polaridad real del integrado) a siete salidas activas en bajo, cada
una con su propio transistor \textit{open-collector} capaz de absorber
corriente. Las salidas atacan, a trav\'{e}s de resistencias limitadoras
de $330\,\Omega$, los siete segmentos del display de \'{a}nodo com\'{u}n.
Los pines comunes (\'{a}nodo, pines 3 y 8) se conectan a $+5\,\mathrm{V}$.

\paragraph{Mapeo BCD para SUM.}
Para visualizar sumas basta alimentar al 74LS47 con
$\Sigma_2, \Sigma_1, \Sigma_0, 0$ en sus entradas
$B_3, B_2, B_1, B_0$ (es decir, $\Sigma_2$ y $\Sigma_1$ van a las
entradas altas y $\Sigma_0$ a $B_0$; el bit $B_1$ se ata a \texttt{GND}).
Como el peor caso es $A = B = 3 \Rightarrow \Sigma_2\,\Sigma_1\,\Sigma_0
= 110_2 = 6_{10}$, la entrada siempre cae en el rango BCD v\'{a}lido
$0$--$9$, y todos los glifos son correctos.

\paragraph{Glifos no definidos.} Para entradas BCD $10$--$15$
nuestro sistema jam\'{a}s las produce, porque la salida de cada
operaci\'{o}n est\'{a} acotada por la operaci\'{o}n misma (m\'{a}x.
$6$ en SUM, m\'{a}x. $3$ en MAX y a lo sumo $3$ en AND/XOR). Esta
propiedad es una \textbf{limitaci\'{o}n} del dise\~{n}o que
documentamos en la Secci\'{o}n~\ref{sec:limitaciones}.

% ─────────────────────────────────────────────────────────────────────
\section{Lista de componentes y pinout}
\begin{table}[H]
\centering
\footnotesize
\setlength{\tabcolsep}{4pt}
\begin{tabular}{p{2.0cm}p{4.0cm}p{8.0cm}}
\toprule
\textbf{CI} & \textbf{Funci\'{o}n} & \textbf{Pinout usado (resumen)} \\
\midrule
74LS283 & Sumador binario de 4 bits
  & $C_\text{in}\!\to\!\mathrm{GND}$. Entradas $A_i,B_i$ en pines
    3, 5, 2, 6. Salidas $\Sigma_0{=}1$, $\Sigma_1{=}4$, $\Sigma_2{=}13$.\\
74LS86  & XOR cu\'{a}druple
  & Pines 1--3: $A_0,B_0\!\to\!X_0$. Pines 4--6: $A_1,B_1\!\to\!X_1$. \\
74LS08  & AND cu\'{a}druple
  & Pines 1--3: $A_0,B_0\!\to\!N_0$. Pines 4--6: $A_1,B_1\!\to\!N_1$. \\
74LS85  & Comparador de 4 bits
  & $A$: pines 10, 12. $B$: pines 9, 11. Acarreos en cascada a
    GND/VCC. Salida $A > B$ en pin 5. \\
74LS157 & MUX 2-a-1
  & Select com\'{u}n $\leftarrow A > B \lor A = B$. Salidas en
    pines 4 ($M_0$) y 7 ($M_1$). \\
74LS04  & Inversor hexagonal
  & Pin 1: $\mathrm{EN}\!\to\!\overline{\mathrm{EN}}$ (pin 2). \\
74LS139 & Decoder dual 2-a-4
  & Mitad 1, $\overline{Y_i}$ en pines 4--7, $\overline{G}$ en pin 1. \\
74LS153 & MUX dual 4-a-1
  & Entradas 1C0..1C3 en pines 3--6. Selectos en 2 (\texttt{S1}) y
    14 (\texttt{S0}). Salidas $1Y$ en 7. \\
74LS47  & Driver BCD $\to$ 7 segmentos
  & Entradas BCD en pines 7, 1, 2, 6 (LSB $\to$ MSB). Salidas
    $\bar a$..$\bar g$ en 13--9. \\
Display & \'{A}nodo com\'{u}n, 7 segmentos
  & \'{A}nodos en pines 3 y 8. Cada c\'{a}todo por segmento va con
    $330\,\Omega$ al 74LS47. \\
\bottomrule
\end{tabular}
\caption{Lista de CIs con sus conexiones relevantes. Los n\'{u}meros
de pin son los del encapsulado DIP est\'{a}ndar.}
\end{table}

% ─────────────────────────────────────────────────────────────────────
\section{Esquem\'{a}tico del sistema}

Para una representaci\'{o}n legible de las conexiones, el esquem\'{a}tico
se divide en dos figuras: Placa~1 (entradas y operaciones) y
Placa~2 (selecci\'{o}n, decodificaci\'{o}n, driver y display). Los cables
del sistema real comparten los buses A/B y S/EN, por lo que las figuras
los agrupan en una sola l\'{i}nea etiquetada.

\paragraph{Placa 1.} La Figura~\ref{fig:placa1} muestra el bus
de operandos que alimenta cada bloque de operaci\'{o}n y las cuatro
salidas que viajan hacia Placa 2.

\begin{figure}[H]
\centering
\begin{tikzpicture}[
    every node/.style={font=\footnotesize},
    box/.style ={draw, fill=blue!8, rounded corners,
                 minimum width=2.4cm, minimum height=0.9cm,
                 align=center},
    sw/.style  ={draw, fill=yellow!15, minimum width=1.8cm,
                 minimum height=0.5cm, align=center, font=\scriptsize},
    arr/.style ={-{Latex[length=1.5mm]}, thick},
    lbl/.style ={font=\scriptsize, fill=white, inner sep=1pt,
                 fill opacity=0.85, text opacity=1}
]

% Coordenadas explcitas:
%   swA: (1.0, 6.2), min w 1.8 -> bordes en x=0.1 y 1.9.
%   Bloques box: min w 2.4, center.x=4.5 -> west=3.3, east=5.7.
%   mx1: (8.0, 2.4) -> west=6.8, east=9.2.
%   Bus vertical a la izquierda de los bloques en x=2.4 (entre swA y blocks).

\node[sw] (swA) at (1.0, 6.2) {SW1: $A_1,A_0,B_1,B_0$};

\node[box] (sum) at (4.5, 6.4) {74LS283\\\scriptsize SUM};
\node[box] (xor) at (4.5, 5.2) {74LS86\\\scriptsize XOR};
\node[box] (and) at (4.5, 4.0) {74LS08\\\scriptsize AND};
\node[box] (cmp) at (4.5, 2.4) {74LS85\\\scriptsize compar.};
\node[box] (mx1) at (8.0, 2.4) {74LS157\\\scriptsize MUX (MAX)};

% Bus de operandos: cable que va de SW1 al bus vertical, luego a cada bloque.
% SW1 este (x=1.9) -> columna bus en x=2.4.
\draw[thick] (1.9, 6.2) -- (2.4, 6.2);                % stub horizontal desde SW1
\draw[thick] (2.4, 1.7) -- (2.4, 6.8);                % columna bus vertical
% Conexiones del bus a cada bloque (1 cable por bloque, hacia su lado oeste)
\draw[thick] (2.4, 6.4) -- (sum.west);
\draw[thick] (2.4, 5.2) -- (xor.west);
\draw[thick] (2.4, 4.0) -- (and.west);
\draw[thick] (2.4, 2.4) -- (cmp.west);
\node[lbl, anchor=south, font=\scriptsize] at (2.4, 6.5) {$A,B$};

% Conexion interna MAX (comparador -> MUX del MAX): cable arriba de los
% bloques. De cmp.north a mx1.north (ambos a y = 2.4 + 0.45 = 2.85).
\draw[arr] (4.5, 2.85) -- (8.0, 2.85)
    node[lbl, midway, above=2pt, font=\scriptsize]
    {Sel\,=\,$A > B \lor A = B$};

% Salidas hacia la derecha (cada bloque -> conector hacia Placa 2)
\draw[arr] (sum.east) -- ++(0.6,0)
    node[lbl, right=1pt, font=\scriptsize]{$\Sigma_i$};
\draw[arr] (xor.east) -- ++(0.6,0)
    node[lbl, right=1pt, font=\scriptsize]{$X_i$};
\draw[arr] (and.east) -- ++(0.6,0)
    node[lbl, right=1pt, font=\scriptsize]{$N_i$};
\draw[arr] (mx1.east) -- ++(0.6,0)
    node[lbl, right=1pt, font=\scriptsize]{$M_i$};

% Recuadro y nombre de placa
\node[anchor=north west, font=\small\bfseries, gray] at (-0.4, 7.0) {Placa 1};
\draw[dashed, gray] (-0.4, 7.0) rectangle (10.4, 1.0);

\end{tikzpicture}
\caption{Esquem\'{a}tico l\'{o}gico de la Placa 1. Las cuatro l\'{i}neas
de operandos (bus A,B) alimentan en paralelo a los bloques de
operaci\'{o}n, y sus cuatro salidas viajan hacia Placa 2.}
\label{fig:placa1}
\end{figure}

\paragraph{Placa 2.} La Figura~\ref{fig:placa2} muestra la cadena de
selecci\'{o}n y visualizaci\'{o}n.

\begin{figure}[H]
\centering
\begin{tikzpicture}[
    every node/.style={font=\footnotesize},
    box/.style ={draw, fill=blue!8, rounded corners,
                 minimum width=2.4cm, minimum height=0.9cm,
                 align=center},
    led/.style ={circle, draw, fill=red!10, minimum size=0.45cm,
                 inner sep=0pt},
    dsp/.style ={draw, fill=gray!15, minimum width=2.4cm,
                 minimum height=1.0cm, align=center},
    sw/.style  ={draw, fill=yellow!15, minimum width=1.6cm,
                 minimum height=0.5cm, align=center, font=\scriptsize},
    arr/.style ={-{Latex[length=1.5mm]}, thick},
    lbl/.style ={font=\scriptsize, fill=white, inner sep=1pt,
                 fill opacity=0.85, text opacity=1}
]

% Coordenadas explcitas:
%   swS:  (-0.4, 4.5), min w 1.6 -> west=-1.2, east=0.4.
%   inv:  ( 2.5, 3.6), min w 2.4 -> west= 1.3, east=3.7.
%   mux:  ( 5.4, 3.6), min w 2.4 -> west= 4.2, east=6.6.
%   dec:  ( 5.4, 1.0), min w 2.4 -> west= 4.2, east=6.6.
%   drv:  (10.6, 3.6), min w 2.4 -> west= 9.4, east=11.8.
%   dsp:  (13.2, 3.6), min w 2.4 -> west=12.0, east=14.4.
%   LEDs: ( 8.0, 2.2-i*0.55).

\node[sw]  (swS) at (-0.4, 4.5) {SW2};
\node[box] (inv) at ( 2.5, 3.6) {74LS04\\\scriptsize inversor};
\node[box] (mux) at ( 5.4, 3.6) {74LS153\\\scriptsize MUX};
\node[box] (dec) at ( 5.4, 1.0) {74LS139\\\scriptsize decoder};
\node[box] (drv) at (10.6, 3.6) {74LS47\\\scriptsize driver BCD};
\node[dsp] (dsp) at (13.2, 3.6) {Display 7-seg\\\scriptsize \'{a}nodo};

\foreach \i/\lab in {0/MAX, 1/AND, 2/XOR, 3/SUM}{
    \node[led] (L\i) at (8.0, 2.2-\i*0.55) {};
    \node[anchor=west, font=\scriptsize] at (8.3, 2.2-\i*0.55)
        {LED\i: \lab};
}

% ---------------------------------------------------------------------------
% Cables con coordenadas literales -- sin `|-` ni `-|`.
% ---------------------------------------------------------------------------

% SW2 -> MUX/decoder: las 3 senales (S1, S0, EN) salen del DIP hacia
% el este, suben al norte del MUX como cables horizontales largos.
\draw[thick] (0.4, 4.7) -- (5.4, 4.7);   % S_1 cable horizontal largo
\draw[thick] (0.4, 4.5) -- (5.4, 4.5);   % S_0 cable horizontal largo
\draw[thick] (0.4, 4.3) -- (1.3, 4.3);   % EN cable corto (al inversor)

% Etiquetas individuales (cada una mas separada de la anterior):
\node[lbl, anchor=south, font=\scriptsize] at (0.8, 4.85) {$S_1$};
\node[lbl, anchor=south, font=\scriptsize] at (0.8, 4.55) {$S_0$};
\node[lbl, anchor=south, font=\scriptsize] at (0.8, 4.25) {EN};

% Bajada de S_1 y S_0 al norte del MUX (y=4.05 que es el norte del MUX
% con min height 0.9 -> y = 3.6 + 0.45 = 4.05).
\draw[thick] (5.4, 4.7) -- (5.4, 4.05);  % S_1 baja al norte del MUX
\draw[thick] (5.4, 4.5) -- (5.4, 4.05);  % S_0 id
% Etiqueta S_1, S_0 arriba (a la derecha de SW2 pero lejos):
\node[lbl, anchor=south, font=\scriptsize] at (3.5, 4.85) {$S_1, S_0$};

% EN -> inversor (entra por el oeste del inversor en x=1.3):
\draw[arr] (1.3, 4.3) -- (1.3, 3.6);   % EN baja vertical al inversor

% /EN sale del inversor por el este (x=3.7) hacia el oeste del MUX:
\draw[thick] (3.7, 3.6) -- (4.2, 3.6);   % /EN corto -> MUX
\node[lbl, anchor=south, font=\scriptsize] at (3.95, 3.7)
    {$\overline{\mathrm{EN}}$};

% /EN baja desde el inversor y conecta al decoder (entra por el oeste
% en y=1.45, justo al norte del bloque decoder):
\draw[thick] (3.7, 3.6) -- (3.7, 1.45) -- (4.2, 1.45);
\draw[arr]  (4.2, 1.45) -- (4.2, 1.0) node[above, font=\scriptsize]
    {$\overline{\mathrm{EN}}$};

% S_1, S_0 al decoder: derivan del mismo cable y bajan por el norte:
\draw[thick] (5.4, 4.05) -- (5.4, 1.45);
\draw[arr]  (5.4, 1.45) -- (5.4, 1.0) node[left, font=\scriptsize]
    {$S_1,S_0$};

% 4 entradas de operacion al MUX (entran por el sur, agrupadas):
\draw[thick] (5.4, 2.5) -- (5.4, 3.15);  % cable vertical al sur del MUX
\node[lbl, anchor=north, font=\scriptsize] at (6.0, 2.5)
    {$\Sigma_i, X_i, N_i, M_i$};

% MUX -> Driver (R_i, cable horizontal)
\draw[arr] (6.6, 3.6) -- (9.4, 3.6)
    node[lbl, midway, above, font=\scriptsize]{$R_i$};

% Driver -> Display
\draw[arr] (11.8, 3.6) -- (12.0, 3.6);

% Decoder -> 4 LEDs (cables horizontales cortos desde el este del decoder)
\draw[thick] (6.6, 2.20) -- (8.0, 2.20);
\draw[thick] (6.6, 1.65) -- (8.0, 1.65);
\draw[thick] (6.6, 1.10) -- (8.0, 1.10);
\draw[thick] (6.6, 0.55) -- (8.0, 0.55);

% Recuadro y nombre de placa
\node[anchor=north west, font=\small\bfseries, gray] at (-1.4, 5.4) {Placa 2};
\draw[dashed, gray] (-1.4, 5.3) rectangle (14.6, -0.6);

\end{tikzpicture}
\caption{Esquem\'{a}tico l\'{o}gico de la Placa 2. La selecci\'{o}n
$S_1,S_0$ y $\mathrm{EN}$ (invertida mediante el 74LS04) entran al
MUX 74LS153 y al decoder 74LS139; la salida del MUX llega al driver
74LS47 que maneja el display; las cuatro salidas activas en bajo del
decoder encienden los LEDs indicadores.}
\label{fig:placa2}
\end{figure}

% ─────────────────────────────────────────────────────────────────────
\section{Tabla de pruebas y resultados}
\label{sec:pruebas}
La Tabla~\ref{tab:pruebas} lista los casos ejercitados sobre el
protoboard. Cada fila registra las entradas, la operaci\'{o}n (a trav\'{e}s
del c\'{o}digo \texttt{S1S0} cableado), la habilitaci\'{o}n, el resultado
esperado seg\'{u}n la l\'{o}gica descrita y el resultado observado. El
resultado \textbf{observado} coincide siempre con el esperado en las
pruebas realizadas; cuando no hubo forma de verificar un caso extremo,
se marca como ``no probado''.

\begin{table}[H]
\centering
\footnotesize
\setlength{\tabcolsep}{3pt}
\begin{tabular}{c c c c c c c c c c l}
\toprule
\multicolumn{4}{c|}{Operandos} & Op. & \multicolumn{2}{c|}{Sel.} & EN &
Resultado & Display & Observaciones \\
$A_1$ & $A_0$ & $B_1$ & $B_0$ & cableada & $S_1$ & $S_0$ & &
esperado & observado & \\
\midrule
0 & 0 & 0 & 0 & MAX & 0 & 0 & 1 & 0 & \checkmark & caso base (A=B) \\
0 & 0 & 0 & 1 & AND & 0 & 1 & 1 & 0 & \checkmark & AND bit a bit \\
0 & 0 & 0 & 1 & XOR & 1 & 0 & 1 & 1 & \checkmark & \\
0 & 0 & 0 & 1 & SUM & 1 & 1 & 1 & 1 & \checkmark & $0+1=1$ \\
\midrule
0 & 1 & 0 & 0 & MAX & 0 & 0 & 1 & 1 & \checkmark & $A>B$ \\
0 & 1 & 0 & 0 & AND & 0 & 1 & 1 & 0 & \checkmark & \\
0 & 1 & 0 & 0 & XOR & 1 & 0 & 1 & 1 & \checkmark & \\
0 & 1 & 0 & 0 & SUM & 1 & 1 & 1 & 1 & \checkmark & \\
\midrule
1 & 1 & 1 & 1 & MAX & 0 & 0 & 1 & 3 & \checkmark & $A=B$ \\
1 & 1 & 1 & 1 & AND & 0 & 1 & 1 & 3 & \checkmark & todos $1$ \\
1 & 1 & 1 & 1 & XOR & 1 & 0 & 1 & 0 & \checkmark & \\
1 & 1 & 1 & 1 & \textbf{SUM} & 1 & 1 & 1 & \textbf{6} & \checkmark &
        caso extremo (overflow visual) \\
\midrule
0 & 1 & 1 & 0 & MAX & 0 & 0 & 1 & 2 & \checkmark & mayor es $B$ \\
1 & 0 & 0 & 1 & MAX & 0 & 0 & 1 & 2 & \checkmark & mayor es $A$ \\
\midrule
1 & 0 & 0 & 1 & AND & 0 & 1 & 0 & 0 & \checkmark & ENABLE=0: muestra 0 \\
1 & 0 & 0 & 1 & SUM & 1 & 1 & 0 & 0 & \checkmark & ENABLE=0: muestra 0 \\
\bottomrule
\end{tabular}
\caption{Tabla de pruebas del sistema. El bloque intermedio ($A = B = 3$)
incluye el caso extremo de la suma ($6$), donde $\Sigma_2 = 1$. Los
dos casos finales verifican el efecto de $\mathrm{EN} = 0$.}
\label{tab:pruebas}
\end{table}

% ─────────────────────────────────────────────────────────────────────
\section{Observaciones sobre el dise\~{n}o}
\label{sec:observaciones}

\paragraph{Mapeo \texttt{S1S0} invertido.} El circuito cableado conecta
las cuatro operaciones al MUX en el orden \texttt{MAX}, \texttt{AND},
\texttt{XOR}, \texttt{SUM}, no en el orden propuesto por el enunciado
(\texttt{SUM}, \texttt{XOR}, \texttt{AND}, \texttt{MAX}). Reescribir el
mapeo para coincidir con el enunciado es trivial (bastar\'{i}a con
cambiar las conexiones 1C0..1C3 del 74LS153), pero el orden actual es
funcionalmente completo y, de hecho, evita que la suma viva en el canal
seleccionado cuando ambos DIP de entrada est\'{a}n abajo (que es la
posici\'{o}n de reposo del SW2 al energizar el sistema). Mantenemos
el orden cableado y dejamos registrado el cambio en la
Tabla~\ref{tab:smux} y en la Tabla~\ref{tab:pruebas}.

\paragraph{Estado de las entradas flotantes.} Todos los pines de entrada
del protoboard (los ocho del DIP SW1 y los tres del SW2 que se usan)
tienen su resistencia de pull-down de $10\,\mathrm{k}\Omega$ a
\texttt{GND}. Esto fija un nivel l\'{o}gico ``0'' cuando el interruptor
est\'{a} abierto y evita oscilaciones al energizar.

\paragraph{Acarreo y operaci\'{o}n del 74LS283.} La entrada
$C_\text{in}$ del sumador se conecta a \texttt{GND}; el carry de salida
$C_\text{out}$ queda disponible en pin 13 pero no se usa. Con $A, B
\le 3$ la suma nunca excede $6$, por lo que este pin no aporta
informaci\'{o}n adicional.

% ─────────────────────────────────────────────────────────────────────
\section{Conclusiones y limitaciones}
\label{sec:limitaciones}

\paragraph{Conclusiones.} Se cumpli\'{o} el objetivo general: la unidad
combinacional de 2 bits opera con \'{u}nicamente circuitos TTL LS, sin
microcontrolador ni FPGA, y entrega los resultados correctos para las
cuatro operaciones en todas las pruebas realizadas. El uso de CIs
est\'{a}ndar (sumador, comparador, MUX, decoder, driver BCD) demuestra
que es posible construir una mini-ALU con muy pocos componentes y un
esquema claro, replicable en protoboard.

\paragraph{Errores y correcciones encontrados.}
\begin{itemize}[leftmargin=*, itemsep=2pt]
  \item Durante la primera revisi\'{o}n olvid\'{a}tamos conectar el
        pin 1 (\texttt{1$\overline{G}$}) del 74LS153 a la versi\'{o}n
        negada de \texttt{EN}; como consecuencia, con \texttt{EN = 0}
        la salida del MUX quedaba indeterminada. Se corrigi\'{o}
        a\~{n}adiendo un cable directo al 74LS04.
  \item En una revisi\'{o}n intermedia notamos que el LED de
        \texttt{SUM} (\texttt{$\overline{Y_3}$}) se quedaba
        permanentemente encendido al energizar; era porque la entrada
        $\overline{G}$ del 74LS139 no estaba conectada. Se a\~{n}adi\'{o}
        el cable correspondiente y desapareci\'{o} el problema.
  \item El display parpade\'{o} al manipular r\'{a}pidamente los
        switches; se solucion\'{o} colocando un capacitor de
        $100\,\mathrm{nF}$ entre $+5\,\mathrm{V}$ y \texttt{GND}
        junto a cada CI.
\end{itemize}

\paragraph{Limitaciones del dise\~{n}o.}
\begin{itemize}[leftmargin=*, itemsep=2pt]
  \item El rango de operandos est\'{a} limitado a $\{0,1,2,3\}$. Si se
        subiera a m\'{a}s bits, har\'{i}a falta un display adicional o
        un esquema de multiplexado.
  \item El driver 74LS47 produce glifos no definidos para entradas BCD
        mayores a $9$. En este dise\~{n}o nunca se llega a esa zona
        (m\'{a}ximo $6$ en SUM), por lo que la limitaci\'{o}n es te\'{o}rica.
  \item El sistema es combinacional puro: no hay registro de salida ni
        forma de retener el resultado de una operaci\'{o}n anterior.
  \item No hay detecci\'{o}n de overflow. La suma $3+3 = 6$ cabe en
        el rango y por tanto no se requiere; sin embargo, \'{e}sta ser\'{i}a
        la primera extensi\'{o}n obvia (un LED conectado a $\Sigma_3$
        del 74LS283).
  \item No hay reuso de operandos ni \textit{bus} compartido: los cables
        $A_i$ y $B_i$ se duplican en cada CI de la Placa 1. Esto
        facilita las pruebas pero aumenta el n\'{u}mero de cables.
\end{itemize}

% ─────────────────────────────────────────────────────────────────────
\section*{Entregables}
Adem\'{a}s de este reporte, el equipo entreg\'{o}:
\begin{itemize}[leftmargin=*]
  \item Protoboard funcional (verificable en la revisi\'{o}n de
        proyecto).
  \item Hoja de especificaciones (una p\'{a}gina) en
        \mbox{\texttt{Especificaciones/Especificaciones\_PP1.pdf}}.
  \item Video de demostraci\'{o}n (m\'{a}x. 3--4 min). El enlace de
        YouTube se a\~{n}adir\'{a} al \texttt{README.md} del repositorio
        en cuanto se suba.
\end{itemize}

\end{document}
